\section{Scalar rigidity for Hermitian flags}
\label{frg-section}

We prove the three assertions which the source combines to
identify the scalar coefficients in its identity-cell calculation.
The conclusion concerns these coefficients and the negative norm;
it does not identify two ordered frames whose vectors can still
differ by lower flag terms. All conventions below agree with the
actual flag and infinitesimal-character definitions.
\cite{alpo-uniform-source}

\subsection{Flags and their infinitesimal characters}

Let $E$ be a number field with an involution $a\mapsto\bar a$,
and let $K\subseteq E$ be a subfield. Let $V$ be a
finite-dimensional $E$-space with a nondegenerate Hermitian form
$h$, conjugate linear in its first argument and linear in its
second. Thus
\[
 h(ax,by)=\bar a b\,h(x,y),\qquad
 h(y,x)=\overline{h(x,y)}.
\]
For $T\in E\setminus\{0\}$, a flag frame of length $r+1$
is a tuple
\[
 F=(e_0,\ldots,e_{r-1},z)
\]
with the $e_i$ linearly independent and
\[
 h(e_i,e_j)=0,\qquad h(e_i,z)=0,\qquad h(z,z)=-T.
\]
Write $f_i=e_i$ for $i<r$, $f_r=z$, and
\begin{gather*}
 X_a=\operatorname{span}_E(e_0,\ldots,e_{a-1})\quad(0\le a\le r),
 \\
 X=X_r,\qquad
 X^\perp=\{y:h(x,y)=0\text{ for every }x\in X\}.
\end{gather*}
For $y\in X$, let $[y]_i$ be its coefficient on $e_i$.
The set of $E$-linear endomorphisms used here is
\begin{equation}\label{frg-lie-definition}
 \mathfrak n_F=\left\{Y:
 \begin{array}{l}
 h(Yx,y)+h(x,Yy)=0\quad(x,y\in V),\\
 YX_a\subseteq X_{a-1}\quad(1\le a\le r),\\
 YX^\perp\subseteq X
 \end{array}\right\}.
\end{equation}
For $Y\in\mathfrak n_F$ define
\begin{equation}\label{frg-character}
 \ell_F(Y)=\sum_{i=0}^{r-1}[Yf_{i+1}]_i,
 \qquad d\psi_F(Y)=\operatorname{Tr}_{E/K}\ell_F(Y).
\end{equation}
Every vector whose coefficient occurs in this formula belongs
to $X$: for $i+1<r$ this follows from
\eqref{frg-lie-definition}, and for $i=r-1$ it follows from
$z\in X^\perp$. When $r=0$, the sum is zero.

Let $F'$ be another such frame, with negative norm parameter
$T'\ne0$. Given scalars $c_0,\ldots,c_r\in E$, the same-flag
condition is exactly
\begin{equation}\label{frg-same-flag}
                         f'_i-c_i f_i\in X_i
                                     \qquad(0\le i\le r).
\end{equation}
The relevance condition needed below is
\begin{equation}\label{frg-relevance}
 d\psi_F(Y)=d\psi_{F'}(Y)
                  \qquad(Y\in\mathfrak n_F\cap\mathfrak n_{F'}).
\end{equation}
It is equality of the actual linear functionals on the
intersection, not just equality on one chosen endomorphism.

The trace pairing has the nondegeneracy needed here without
an additional hypothesis. Indeed $E/K$ is a finite extension
of characteristic zero. If $b\ne0$, then
\[
 \operatorname{Tr}_{E/K}(b\,b^{-1})=[E:K]\ne0.
\]
Consequently
\begin{equation}\label{frg-trace-nondegenerate}
 \bigl(\operatorname{Tr}_{E/K}(ba)=0\text{ for all }a\in E\bigr)
                                      \quad\Longrightarrow\quad b=0.
\end{equation}
This proves the source's explicit trace-pairing assumption in
the present number-field setting.

\subsection{Relevance forces the flag scalars to agree}

\begin{lemma}\label{frg-triangular}
Under \eqref{frg-same-flag}, every $c_i$ is nonzero and
$X'_a=X_a$ for $0\le a\le r$.
\end{lemma}
\begin{proof}
The vectors $f_0,\ldots,f_r$ are independent. For if
$\sum_{i<r}a_i e_i+a_rz=0$, pairing with $z$ in the first
argument gives $-a_rT=0$. Then $a_r=0$, and independence
of the $e_i$ gives all other coefficients zero. The same
argument applies to $F'$.
Equation \eqref{frg-same-flag} writes the independent tuple
$F'$ as a triangular change of the independent tuple $F$,
with diagonal entries $c_i$. Its matrix is invertible,
so every $c_i\ne0$. Each initial triangular block is likewise
invertible, giving $X'_a=X_a$.
\end{proof}

\begin{proposition}\label{frg-relevant-scalars}
If \eqref{frg-same-flag} and \eqref{frg-relevance} hold, then
\begin{equation}\label{frg-all-scalars}
                              c_i=c_r\qquad(0\le i\le r).
\end{equation}
\end{proposition}
\begin{proof}
There is nothing to prove for $r=0$. Fix $0\le i<r$ and
$a\in E$. Nondegeneracy of $h$ and independence of the $f_j$
give a vector $v\in V$ such that
\begin{equation}\label{frg-dual-vector}
                     h(v,f_j)=a\,\delta_{j,i+1}
                                        \qquad(0\le j\le r).
\end{equation}
For completeness, extend the $f_j$ to a basis $b_1,\ldots,b_d$.
The Gram matrix $G=(h(b_j,b_k))$ is invertible, since a vector
orthogonal to the whole basis would contradict nondegeneracy.
Prescribe the desired values $h(v,b_k)$, using zero on the
additional basis vectors. If $v=\sum_j t_jb_j$, these equations
are $\bar t^{\,t}G=(h(v,b_k))_k$. Invert $G$ and apply the
involution to solve for every $t_j$. This proves
\eqref{frg-dual-vector}.

Define the $E$-linear endomorphism
\begin{equation}\label{frg-rank-two}
                       Y(y)=h(v,y)e_i-h(e_i,y)v.
\end{equation}
It is skew Hermitian. Indeed the two terms of $h(Yx,y)$ are
$h(x,v)h(e_i,y)-h(x,e_i)h(v,y)$, which cancel the two terms
of $h(x,Yy)$.
Also $h(v,X_{i+1})=0$ by \eqref{frg-dual-vector}.
For $x\in X_b$ we have $h(e_i,x)=0$. If $b\le i+1$, then
$Yx=0$; if $b\ge i+2$, then $Yx$ is a multiple of
$e_i\in X_{b-1}$. Finally, for $x\in X^\perp$ the second
term in \eqref{frg-rank-two} vanishes, and $Yx\in X$.
Thus $Y\in\mathfrak n_F$.

By Lemma~\ref{frg-triangular}, all the spaces $X'_b$ equal
$X_b$. The same argument, using $e_i\in X'_{i+1}$ and
$h(v,X'_{i+1})=0$, proves $Y\in\mathfrak n_{F'}$.
As $h(e_i,f_j)=0$ for every $j$, equation
\eqref{frg-dual-vector} gives immediately
\begin{equation}\label{frg-ell-unprimed}
                                  \ell_F(Y)=a.
\end{equation}

We compute the primed functional, retaining all triangular
terms. Since $f'_j\in X+Ez$ and $h(e_i,X+Ez)=0$,
\[
                            Yf'_{j+1}=h(v,f'_{j+1})e_i.
\]
If $j<i$, then $f'_{j+1}\in X_{i+1}$ and the scalar on
the right is zero. If $j>i$, then $e_i\in X'_{i+1}$
has zero coefficient on $e'_j$. For $j=i$, triangularity
gives
\[
 [e_i]'_i=c_i^{-1},\qquad
 h(v,f'_{i+1})=c_{i+1}a,
\]
because the error in $f'_{i+1}-c_{i+1}f_{i+1}$ lies in
$X_{i+1}$, which is orthogonal to $v$ in the required order.
Only this term survives. Hence
\begin{equation}\label{frg-ell-primed}
                         \ell_{F'}(Y)=\frac{c_{i+1}}{c_i}a.
\end{equation}
Relevance and \eqref{frg-ell-unprimed}--\eqref{frg-ell-primed}
give
\[
 \operatorname{Tr}_{E/K}\left(\left(\frac{c_{i+1}}{c_i}-1\right)a\right)=0
                                      \qquad(a\in E).
\]
Equation \eqref{frg-trace-nondegenerate} implies
$c_{i+1}=c_i$. This holds for each $i<r$, proving
\eqref{frg-all-scalars}.
\end{proof}

\subsection{An integral unit and the negative norm}

Assume now that the involution is complex conjugation at every
embedding:
\begin{equation}\label{frg-all-embeddings}
                 \sigma(\bar b)=\overline{\sigma(b)}
                    \quad(b\in E,\ \sigma:E\hookrightarrow\mathbb C).
\end{equation}
This includes the actual CM-field involution. Let $T,T'$ be
positive integers and suppose the final scalar is the image
of a unit $\zeta\in\mathcal O_E^\times$.

\begin{proposition}\label{frg-unit-norm-root}
Under \eqref{frg-same-flag} and \eqref{frg-all-embeddings},
if $c_r=\zeta\in\mathcal O_E^\times$, then $T=T'$ and
$\zeta$ is a root of unity. Relevance is not needed for this
assertion.
\end{proposition}
\begin{proof}
The final equation in \eqref{frg-same-flag} says
$z'=\zeta z+x$ with $x\in X$. The subspace $X$ is
totally isotropic and orthogonal to $z$, so
\[
                 -T'=h(z',z')=-\bar\zeta\zeta T,
 \qquad T'=\bar\zeta\zeta T.
\]
Let $d=[E:\mathbb Q]$. The norm of an algebraic unit is
$\pm1$: multiplication by the unit is an automorphism of
the free abelian group $\mathcal O_E$, so its determinant
is $\pm1$. The involution is a $\mathbb Q$-automorphism,
and conjugates the multiplication operators of $\zeta$
and $\bar\zeta$. Thus their norms agree. Taking norms in
the last display gives
\[
                 (T')^d=N_{E/\mathbb Q}(\zeta)^2T^d=T^d.
\]
Positivity gives $T'=T$, and cancellation gives
$\bar\zeta\zeta=1$. Equation \eqref{frg-all-embeddings}
now implies $|\sigma(\zeta)|=1$ for every embedding.

Here is the elementary finite-polynomial argument completing
the root-of-unity assertion. For every integer $m\ge1$,
the characteristic polynomial of multiplication by $\zeta^m$
on $\mathcal O_E$ is
\[
                P_m(t)=\prod_{\sigma:E\hookrightarrow\mathbb C}
                                  (t-\sigma(\zeta)^m)\in\mathbb Z[t].
\]
Its coefficient of $t^{d-j}$ has absolute value at most
$\binom dj$, since it is a sum of that many products of
complex numbers of absolute value one. Only finitely many
monic integer polynomials of degree $d$ have these coefficient
bounds. Their union of root sets is finite. In a fixed
embedding the powers $\sigma(\zeta)^m$ therefore take only
finitely many values. Two positive powers are equal;
injectivity and $\zeta\ne0$ give $\zeta^n=1$ for some
$n>0$.
\end{proof}

\subsection{Reduction at an odd unramified prime}

\begin{lemma}\label{frg-root-congruence}
Let $P$ be a nonzero prime ideal of $\mathcal O_E$, and
let $p\ge3$ be a rational prime such that
\begin{equation}\label{frg-prime-condition}
                         p\in P\quad\text{and}\quad p\notin P^2.
\end{equation}
If $\zeta\in\mathcal O_E$ is a root of unity and
$\zeta-1\in P$, then $\zeta=1$.
\end{lemma}
\begin{proof}
Suppose $\zeta\ne1$, and let its order be $m>1$.
Choose a prime $q\mid m$ and put $\eta=\zeta^{m/q}$.
Then $\eta$ has order exactly $q$, and $\eta-1\in P$.
The roots $1,\eta,\ldots,\eta^{q-1}$ are distinct, so
factoring $t^q-1$ over $E$, dividing by $t-1$, and
evaluating at $t=1$ gives
\begin{equation}\label{frg-prime-product}
                     q=\prod_{j=1}^{q-1}(1-\eta^j)\in P^{q-1}.
\end{equation}
Indeed each factor lies in $P$, because $\eta\equiv1\pmod P$.
In particular $q\in P$. If $p\ne q$, integers $a,b$
with $ap+bq=1$ would imply $1\in P$, a contradiction.
Thus $q=p$. Since $p-1\ge2$, equation
\eqref{frg-prime-product} implies $p\in P^2$, contrary
to \eqref{frg-prime-condition}.
\end{proof}

\begin{theorem}[The actual flag-scalar conclusion]
\label{frg-identity-cell}
Let $E,V,h,K,F,F',c$ satisfy the definitions above, with
$T,T'$ positive integers. Suppose
\eqref{frg-same-flag}, \eqref{frg-relevance} and
\eqref{frg-all-embeddings} hold. Suppose further that
$c_r=\zeta\in\mathcal O_E^\times$ and that, for a rational
prime $p\ge3$ and a prime ideal $P$,
\[
                         p\in P\setminus P^2,
                         \qquad\zeta\equiv1\pmod P.
\]
Then
\begin{equation}\label{frg-identity-conclusion}
                         T=T',\qquad c_0=\cdots=c_r=1.
\end{equation}
\end{theorem}
\begin{proof}
Proposition~\ref{frg-unit-norm-root} gives $T=T'$ and
makes $\zeta$ a root of unity. Lemma~\ref{frg-root-congruence}
gives $\zeta=1$. Finally Proposition~\ref{frg-relevant-scalars}
gives $c_i=c_r=\zeta$ for every $i$.
\end{proof}

The restrictions in the theorem cannot simply be discarded.
At the prime $2$, the unit $-1\in\mathbb Z$ satisfies
$-1\equiv1\pmod2$ and $2\notin(2)^2$, but is not $1$.
For ramified odd primes, take a primitive cube root $\omega$
in $E=\mathbb Q(\omega)$ and $P=(1-\omega)$.
Then $\omega\equiv1\pmod P$ and
$3=-\omega^{-1}(1-\omega)^2\in P^2$, so the claimed
root conclusion would again fail without the exclusion of
$P^2$. The unit condition in Proposition~\ref{frg-unit-norm-root}
is also necessary: for $V=\mathbb Q(i)$ with
$h(x,y)=-\bar x y$, take $r=0$, $z=1$, $z'=2$.
The same-flag relation has scalar $2$, but $T=1$ and $T'=4$.

Even the full theorem does not assert $F=F'$ as ordered
tuples. When $r\ge1$, replacing $z$ by $z+e_0$ preserves
its norm, all flag conditions and all scalars $c_i=1$.
It also preserves the infinitesimal character, since
$Ye_0=0$ for $Y\in\mathfrak n_F$, and the spaces
\eqref{frg-lie-definition} depend only on the $e_i$.
The source's precise conclusion is \eqref{frg-identity-conclusion},
which is what has been proved. No counterexample satisfying
those actual hypotheses was found.


% Ordinary mathematics for the vanishing of degenerate off-diagonal terms.
% Labels in this fragment have the prefix df-.
